\documentclass[11pt]{article}
\usepackage{amsmath,amssymb,amsthm}
\usepackage[margin=1in]{geometry}
\usepackage[T1]{fontenc}
\newcommand{\C}{\mathbb{C}}
\newcommand{\Z}{\mathbb{Z}}
\newcommand{\Q}{\mathbb{Q}}
\newcommand{\HH}{\mathbb{H}}
\newcommand{\RR}{\mathcal{R}}
\newcommand{\GG}{\mathcal{G}}
\newcommand{\FF}{\mathcal{F}}
\newcommand{\JJ}{\mathcal{J}}
\theoremstyle{definition}
\newtheorem{defn}{Definition}
\newtheorem{keyformula}{Key Formula}


\begin{document}

The companion paper (Alfes--Ono, \emph{Higher Dyson ranks and meromorphic elliptic
corrections}) develops, for each $m\ge 2$ and $d:=2m+1$, a finite elliptic correction theory
for the higher Dyson series $\RR_m$. The overall architecture, meromorphic Jacobi forms,
Appell--Lerch kernels, weak-Maass and mock-modular background, the Deuring/CM and modular
input---is either classical or is the kind of analytic bookkeeping (normal convergence on
compacta, holomorphy of coefficients) that we take as established. What is genuinely new, and
where a hidden sign, exponent, root-of-unity phase, or normalization error could enter
unnoticed, is a short chain of \emph{algebraic} identities in $q$-series and rational functions.
We isolate exactly those identities here and target them for machine verification. We do not
attempt to formalize the paper as a whole.

Two features make these identities formalizable. First, each is an identity of \emph{formal}
objects: an equality in $\C(u)$, in a Laurent ring $\C((q))[u]$, or---after evaluating the
elliptic variable at a fixed $\tau$---an equality of complex numbers. No analytic limit enters
the statements themselves. Second, the ``modular'' or ``$q$-series'' inputs each identity needs
can be introduced as opaque symbols subject only to the finitely many algebraic relations
actually used (the two product $q$-shift laws, the theta triple-product shifts). We therefore
introduce those symbols and relations directly, and phrase every Key Formula as a
self-contained algebraic equality.

Throughout, $m\ge 2$ is an integer and $d:=2m+1$ (so $d\ge 5$ is odd and $d-3=2m-2$ is even and
$\ge 2$). We fix $x\in\C^\times$, and when a root of unity is needed we write $\omega:=e^{2\pi i/d}$.
The tags in brackets name the corresponding statement in the paper.

%==============================================================================
\section*{Formalization conventions}
%==============================================================================

These conventions are part of the specification: they fix how the objects are to be modeled,
and any formalization that departs from them is proving a different statement. They do not
weaken any Key Formula.

\begin{enumerate}
\item[(C1)] \textbf{Opaque symbols, not series.} The objects
$G_d,\,M_d,\,P_d,\,X_j^{(m)},\,\FF_x,\,\GG_{d,x},\,K_{d,x},\,\Phi_{d,x},\,\mathcal H_{d,x}$
and the correction kernels
$\Psi_d^{[x]},\,\Omega_{d,r}^{[x]},\,\Omega_{d,\mathrm{ext}}^{[x]},\,
\mathcal B_d^{[x]},\,\mathcal A_{d,r}^{[x]},\,\mathcal A_{d,\mathrm{ext}}^{[x]}$
are \emph{opaque parameters} (functions to $\C$, with arguments in $\C$ or $\C^\times$).
They must \emph{not} be defined by infinite sums, products, or \texttt{tsum}. Their only
permitted properties are the finitely many algebraic relations written for them below---the
$G_d$ shift, the $M_d$ and $P_d$ multiplicative $q$-shift laws, the kernel $q$-shift and
$\tau$-shift laws, and the $\vartheta_1$ triple-product normalization and quasi-periodicities.
Each such relation is introduced as a \emph{hypothesis} of the statement that uses it. No
summability or convergence claim is ever a proof obligation; where a kernel is manipulated,
its shift law is the hypothesis that licenses it.

\item[(C2)] \textbf{Units.} Every symbol that is inverted or raised to a negative power is a
unit: take $x\in\C^\times$ (equivalently $x\ne0$) and $u\in\C^\times$ throughout. In
particular the hypothesis ``$x^d\ne1$'' does not license $x^{-1}$ on its own; carry $x\ne0$
(or $x\in\C^\times$) explicitly wherever $x^{-1}$ or $x^{-d}$ appears.

\item[(C3)] \textbf{Global setup carried everywhere.} Fix once $d:=2m+1$ with $m\ge2$ an
integer, and phrase every statement in terms of $m$ (so $d\ge5$ odd, $d-3=2m-2\ge2$ even).
This hypothesis travels with every Key Formula.

\item[(C4)] \textbf{Pinned constants (do not re-derive).} The following normalizations are
fixed data, to be taken as named givens rather than recomputed:
\[
\text{$1/u$ partial-fraction coefficient}=-X_0^{(m)}(x;q)=T_{d,x}(0;q)=R_{d,x}(0;q);
\]
\[
\vartheta_1(\tau,z+\tau)=-q^{-1/2}u^{-1}\vartheta_1(\tau,z)
\quad(\text{half-integer power }-\tfrac12,\ \text{not }-1);
\]
\[
\vartheta_1(\tau,z)=iq^{1/8}u^{-1/2}J(u;q),
\qquad
\eta(\tau)=q^{1/24}(q;q)_\infty
\ \Rightarrow\ \eta(\tau)^{24}=q\,(q;q)_\infty^{24};
\]
and in the level-$5$ carrier the scalar prefactor is $q^{1/4}$ (from
$q^{(m+1)/12}=q^{3/12}=q^{1/4}$ at $m=2$).

\item[(C5)] \textbf{Rational-function reading.} Identities among rational functions of $u$
(Key Formulas~\ref{KF:functional},~\ref{KF:defect},~\ref{KF:residues},~\ref{KF:kernelshift})
are equalities in the field $\C(u)$ (equivalently, equalities of numerators after clearing
the common denominator), not pointwise identities carrying side conditions. Statements in the
elliptic variable (Key Formulas~\ref{KF:Pshift}--\ref{KF:carrier}) are equalities of the
opaque functions under the stated shifts.

\item[(C6)] \textbf{Statements are conclusions, not hypotheses.} Each Key Formula is the thing
to be \emph{proved}. In particular Key Formula~\ref{KF:functional} is an existential:
$\exists\,R_{d,x}\in q^{-1}\C[[q]][u]$ with $\deg_u R_{d,x}\le d-2$, $R_{d,x}(0;q)=-X_0^{(m)}(x;q)$,
and the displayed functional equation. The recurrence, the constant-term value, the twist
laws, and the decay/convergence facts may appear as \emph{hypotheses only} when the
convention above designates them as opaque-symbol shift laws; they may never be assumed as
stand-ins for the conclusion being proved.
\end{enumerate}

%==============================================================================
\section*{Ground data and definitions}
%==============================================================================

\begin{defn}[The building-block quotient $G_d$ and its shift]\label{def:G}
Work in the formal Laurent field $\C((q))$ with a marked unit $x\in\C^\times$. The object
$G_d(x;q):=(x^dq^d;q^d)_\infty/(xq;q)_\infty$ enters only through the single algebraic relation
it satisfies, which we take as its defining property:
\[
\frac{G_d(x;q)}{G_d(xq;q)}=\frac{1-x^dq^d}{1-xq}
=\sum_{r=0}^{d-1}x^rq^r .
\]
Equivalently $G_d(x;q)=\bigl(\sum_{r=0}^{d-1}x^rq^r\bigr)\,G_d(xq;q)$, and more generally, for
$n\ge0$,
\[
G_d(xq^n;q)=\Bigl(\sum_{r=0}^{d-1}x^rq^{r(n+1)}\Bigr)G_d(xq^{n+1};q).
\]
\end{defn}

\begin{defn}[Deformed higher Dyson values and their generating series]\label{def:X}
For $j\in\Z$ set
\[
X_j:=X_j^{(m)}(x;q):=\sum_{n\ge0}q^{n^2+jn}G_d(xq^n;q).
\]
\emph{Power-series constraint (essential).} Because the sum runs over $n\ge0$ with exponents
$n^2+jn$ bounded below, and $G_d(xq^n;q)$ is itself a power series in $q$, each $X_j$ is a
genuine \emph{power series} in $q$: $X_j\in\C[[q]]$ (nonnegative powers of $q$ only), and in
particular $X_0\in\C[[q]]$. This is a defining property of the $X_j$, \textbf{not} a derived
one: the objects must be typed as elements of $\C[[q]]$ (equivalently, carry the hypothesis
$X_j\in\C[[q]]$ for every $j$ used), \emph{not} as arbitrary two-sided elements of $\C((q))$.
Every negative power of $q$ in the construction originates from the single term $q^{-1}X_0$ in
the boundary closure (Key Formula~\ref{KF:boundary}), which is why $R_{d,x}\in q^{-1}\C[[q]][u]$
carries exactly one order of pole in $q$; without the constraint $X_0\in\C[[q]]$ the functional
equation (Key Formula~\ref{KF:functional}) is \emph{false} (e.g.\ $X_0=q^{-1}$ forces
$q\cdot X_0=q^{-1}\notin\C[[q]]$, so no valid $R_{d,x}$ exists). Then
\[
\FF_x(t):=\sum_{j\ge0}X_j\,t^j,
\qquad
\GG(u):=\GG_{d,x}(u;q):=x^{-2}\FF_x(xu).
\]
Write $G:=G_d(x;q)$ and $C_x:=\sum_{r=0}^{d-1}x^r=\dfrac{1-x^d}{1-x}$ (for $x\ne1$).
\end{defn}

\begin{defn}[The two normalizing products and their $q$-shifts]\label{def:MP}
Let $J(u;q):=(u;q)_\infty(q/u;q)_\infty(q;q)_\infty$ be the Jacobi triple product. Set
\[
M_d(u;q):=\frac{(u^d;q^d)_\infty}{(u;q)_\infty\,J(u;q)^{\,d-3}},
\qquad
P_d(u;q):=u^{-m-1}J(-u;q)^{\,d+1}\frac{(q/u;q)_\infty}{(q^d/u^d;q^d)_\infty}.
\]
These enter the algebra \emph{only} through their multiplicative $q$-shift laws, which we take
as defining relations:
\[
M_d(qu;q)=\frac{u^{d-3}(1-u)}{1-u^d}\,M_d(u;q),
\qquad
P_d(qu;q)=q^{-m-1}u^{-2}\,\frac{1-u}{1-u^d}\,P_d(u;q).
\]
The normalized generating series is $K_{d,x}(u;q):=\dfrac{M_d(u;q)}{u}\,\GG_{d,x}(u;q)$.
\end{defn}

\begin{defn}[The defect polynomial and its clearing]\label{def:RT}
$R_{d,x}(u;q)\in q^{-1}\C[[q]][u]$ is the polynomial of degree $\le d-2$ produced in
Key Formula~\ref{KF:functional}; define
\[
T_{d,x}(u;q):=(1-xu)(1-u)R_{d,x}(u;q)-C_x\,u^{d-1}(1-u)\,G_d(x;q),
\]
a polynomial in $u$ of degree $\le d$.
\end{defn}

\begin{defn}[Correction kernels and Appell-type sums]\label{def:kernels}
For $u=e^{2\pi i w}$ set (rational-function level)
\[
\Psi_d^{[x]}(u):=\sum_{n\ge0}x^{-2n-2}\frac{q^{-n}M_d(q^nu;q)}{u},
\quad
\Omega_{d,r}^{[x]}(u):=\sum_{n\ge0}\frac{x^{-2n-2}M_d(q^nu;q)}{1-\omega^rq^nu},
\quad
\Omega_{d,\mathrm{ext}}^{[x]}(u):=\sum_{n\ge0}\frac{x^{-2n-2}M_d(q^nu;q)}{1-xq^nu},
\]
and (elliptic-variable level, with $u=e^{2\pi i w}$)
\[
\mathcal B_d^{[x]}:=\sum_{n\ge0}x^{-2n-2}q^{mn^2}u^{2mn}P_d(q^nu;q),
\qquad
\mathcal A_{d,r}^{[x]}:=\sum_{n\ge0}x^{-2n-2}q^{mn^2+n}u^{2mn+1}\frac{P_d(q^nu;q)}{1-\omega^rq^nu},
\]
and $\mathcal A_{d,\mathrm{ext}}^{[x]}$ likewise with $1-\omega^rq^nu$ replaced by $1-xq^nu$.
(Only the $q$-shift laws below are needed; convergence is assumed.)
\end{defn}

\begin{defn}[Theta data for the carrier]\label{def:theta}
Let $\eta,\vartheta_1$ be Dedekind's eta and the odd Jacobi theta function. On the relevant
congruence subgroup $\Gamma(24d^2)$ they enter only through the quasi-periodicity relations
\[
\vartheta_1(\tau,z+1)=-\vartheta_1(\tau,z),
\qquad
\vartheta_1(\tau,z+\tau)=-q^{-1/2}u^{-1}\vartheta_1(\tau,z),
\]
(and the $d$-scaled versions with $\tau\mapsto d\tau$, $z\mapsto dz$, $u\mapsto u^d$, $q\mapsto q^d$),
together with the triple-product normalization $\vartheta_1(\tau,z)=iq^{1/8}u^{-1/2}J(u;q)$.
The Jacobi carrier is
\[
\JJ_{m,d}(\tau,z):=\eta(\tau)^{\,2-4m}\,\frac{\vartheta_1(\tau,z)^{\,4m+1}}{\vartheta_1(d\tau,dz)}.
\]
\end{defn}

%==============================================================================
\section*{The key identities}
%==============================================================================

\begin{keyformula}[Root-of-unity collapse of the inhomogeneity]\label{KF:collapse}
With $C_x=\sum_{r=0}^{d-1}x^r$, if $x^d=1$ but $x\ne1$ then
\[
C_x=\frac{1-x^d}{1-x}=0 .
\]
Consequently the inhomogeneous term $C_x\,G_d(x;q)$ of the recurrence
(Key Formula~\ref{KF:recurrence}) vanishes at every nontrivial $d$-th root of unity.
\end{keyformula}

\begin{keyformula}[The fundamental $(2m+1)$-term recurrence; Prop.~2.2, eq.~(2.1)]\label{KF:recurrence}
For every $j\in\Z$,
\[
\sum_{r=0}^{d-1}x^r\,X_{j+r-2}^{(m)}(x;q)\;-\;q^{\,j-1}X_j^{(m)}(x;q)
\;=\;\Bigl(\sum_{r=0}^{d-1}x^r\Bigr)G_d(x;q).
\]
The verification is the exponent identity
\[
n^2+(j+r)n+r=(n+1)^2+(j+r-2)(n+1)+(1-j),
\]
combined with the shift $G_d(xq^n;q)=\bigl(\sum_r x^rq^{r(n+1)}\bigr)G_d(xq^{n+1};q)$ of
Definition~\ref{def:G} and reindexing $N=n+1$, where the missing $N=0$ term contributes
exactly $-G_d(x;q)$ per summand.
\end{keyformula}

\begin{keyformula}[The two negative-index boundary closures; eqs.~(4.5),(4.6)]\label{KF:boundary}
Specializing Key Formula~\ref{KF:recurrence} at $j=1$ and $j=0$ and eliminating gives
\[
X_{-1}=C_xG-\sum_{r=1}^{d-1}x^rX_{r-1}+X_1,
\qquad
X_{-2}=q^{-1}X_0-xX_1+x^dX_{d-2}+(1-x)C_xG .
\]
These are the only closures needed to sum the recurrence over $j\ge0$.
\end{keyformula}

\begin{keyformula}[The master functional equation and the constant term of $R_{d,x}$; Thm.~4.1, eqs.~(1.8),(4.1)]\label{KF:functional}
Assume the power-series constraint of Definition~\ref{def:X}: $X_j^{(m)}(x;q)\in\C[[q]]$ for
every $j$ (in particular $X_0^{(m)}(x;q)\in\C[[q]]$). This hypothesis is required---without it
the statement is false. Then there is a unique $R_{d,x}(u;q)\in q^{-1}\C[[q]][u]$ of degree
$\le d-2$ with
\[
\GG_{d,x}(qu;q)
=x^2q\,\frac{u^d-1}{u^{d-3}(u-1)}\,\GG_{d,x}(u;q)
+\frac{q}{u^{d-3}}\,R_{d,x}(u;q)
-C_x\,q\,\frac{u^2}{1-xu}\,G_d(x;q),
\]
and moreover $R_{d,x}(0;q)=-X_0^{(m)}(x;q)$. The two rational-function reductions that produce
this form are
\[
\sum_{r=0}^{d-1}u^{2-r}=\frac{u^d-1}{u^{d-3}(u-1)},
\qquad
1+xu-\frac{1}{1-xu}=-\frac{x^2u^2}{1-xu}.
\]
\end{keyformula}

\begin{keyformula}[$M_d$-clearing to the twisted defect equation; eqs.~(4.10),(4.3)]\label{KF:defect}
Using the $M_d$ shift of Definition~\ref{def:MP} to multiply the master equation by
$M_d(qu;q)/(qu)$, the homogeneous coefficient collapses,
\[
\frac{M_d(qu;q)}{u}\cdot x^2q\,\frac{u^d-1}{u^{d-3}(u-1)}=x^2\,\frac{M_d(u;q)}{u},
\]
and the whole identity becomes, with $T_{d,x}$ as in Definition~\ref{def:RT},
\[
K_{d,x}(qu;q)-x^2K_{d,x}(u;q)
=\frac{M_d(u;q)}{u(1-u^d)(1-xu)}\,T_{d,x}(u;q).
\]
\end{keyformula}

\begin{keyformula}[Explicit partial-fraction residues; Thm.~4.3, eqs.~(4.18),(4.19)]\label{KF:residues}
Since $x^d\ne1$, the $d+2$ linear factors $u,\,1-\omega^ru\ (0\le r\le d-1),\,1-xu$ are distinct, so
\[
\frac{T_{d,x}(u;q)}{u(1-u^d)(1-xu)}
=-\frac{X_0^{(m)}(x;q)}{u}
+\sum_{r=0}^{d-1}\frac{D_{r,x}(q)}{1-\omega^ru}
+\frac{E_x(q)}{1-xu},
\]
with the residues determined by evaluation at the poles:
\[
-X_0^{(m)}(x;q)=\bigl.\tfrac{T_{d,x}(u;q)}{(1-u^d)(1-xu)}\bigr|_{u=0}=T_{d,x}(0;q),
\quad
D_{r,x}(q)=\frac{\omega^r}{d\,(1-x\omega^{-r})}\,T_{d,x}(\omega^{-r};q),
\quad
E_x(q)=\frac{x}{1-x^{-d}}\,T_{d,x}(x^{-1};q).
\]
\end{keyformula}

\begin{keyformula}[Kernel $q$-shift identities matching the defect; Prop.~4.2, eqs.~(4.14)--(4.16)]\label{KF:kernelshift}
The three correction kernels of Definition~\ref{def:kernels} satisfy, by termwise index shift,
\[
\Psi_d^{[x]}(qu)-x^2\Psi_d^{[x]}(u)=-\frac{M_d(u;q)}{u},
\quad
\Omega_{d,r}^{[x]}(qu)-x^2\Omega_{d,r}^{[x]}(u)=-\frac{M_d(u;q)}{1-\omega^ru},
\quad
\Omega_{d,\mathrm{ext}}^{[x]}(qu)-x^2\Omega_{d,\mathrm{ext}}^{[x]}(u)=-\frac{M_d(u;q)}{1-xu}.
\]
Combined with Key Formulas~\ref{KF:defect} and~\ref{KF:residues} these force
\[
K_{d,x}(u;q)=X_0^{(m)}(x;q)\,\Psi_d^{[x]}(u)
-\sum_{r=0}^{d-1}D_{r,x}(q)\,\Omega_{d,r}^{[x]}(u)
-E_x(q)\,\Omega_{d,\mathrm{ext}}^{[x]}(u).
\]
\end{keyformula}

\begin{keyformula}[$P_d$-shift and the index-$m$ elliptic multiplier; Prop.~4.4, eqs.~(4.22),(4.21)]\label{KF:Pshift}
Using the $P_d$ shift of Definition~\ref{def:MP} and $d=2m+1$, the homogeneous coefficient of
the $P_d$-normalized equation collapses to the standard index-$m$ multiplier:
\[
P_d(qu;q)\cdot x^2q\,\frac{u^d-1}{u^{d-3}(u-1)}=x^2\,q^{-m}u^{-2m}\,P_d(u;q).
\]
Consequently $\Phi_{d,x}:=P_d\,\GG_{d,x}$ satisfies $\Phi_{d,x}(\tau,w+1)=\Phi_{d,x}(\tau,w)$ and
\[
\Phi_{d,x}(\tau,w+\tau)=x^2q^{-m}u^{-2m}\Phi_{d,x}(\tau,w)
+q^{-m}u^{-2m}P_d(u;q)\!\left(-X_0^{(m)}(x;q)
+\sum_{r=0}^{d-1}\frac{D_{r,x}(q)\,u}{1-\omega^ru}
+\frac{E_x(q)\,u}{1-xu}\right).
\]
\end{keyformula}

\begin{keyformula}[Appell-kernel $\tau$-shift laws that cancel the defect; Prop.~4.5]\label{KF:appellshift}
The elliptic-variable kernels of Definition~\ref{def:kernels} satisfy, via
$mn^2+2mn=m(n+1)^2-m$ and $mn^2+n+2mn+1=m(n+1)^2+(n+1)-m$,
\[
\mathcal B_d^{[x]}(\tau,w+\tau)=x^2q^{-m}u^{-2m}\mathcal B_d^{[x]}(\tau,w)-q^{-m}u^{-2m}P_d(u;q),
\]
\[
\mathcal A_{d,r}^{[x]}(\tau,w+\tau)=x^2q^{-m}u^{-2m}\mathcal A_{d,r}^{[x]}(\tau,w)
-q^{-m}u^{-2m}\frac{uP_d(u;q)}{1-\omega^ru},
\]
and the analogous law for $\mathcal A_{d,\mathrm{ext}}^{[x]}$ with $1-\omega^ru\to1-xu$.
These principal-part defects match those in Key Formula~\ref{KF:Pshift} term for term.
\end{keyformula}

\begin{keyformula}[The corrected function is elliptic; Cor.~4.6, eqs.~(4.26),(4.27)]\label{KF:corrected}
Define
\[
\mathcal H_{d,x}(\tau,w):=\Phi_{d,x}(\tau,w)
-X_0^{(m)}(x;q)\,\mathcal B_d^{[x]}(\tau,w)
+\sum_{r=0}^{d-1}D_{r,x}(q)\,\mathcal A_{d,r}^{[x]}(\tau,w)
+E_x(q)\,\mathcal A_{d,\mathrm{ext}}^{[x]}(\tau,w).
\]
Then the defects of Key Formulas~\ref{KF:Pshift} and~\ref{KF:appellshift} cancel and
\[
\mathcal H_{d,x}(\tau,w+1)=\mathcal H_{d,x}(\tau,w),
\qquad
\mathcal H_{d,x}(\tau,w+\tau)=x^2q^{-m}u^{-2m}\,\mathcal H_{d,x}(\tau,w).
\]
\end{keyformula}

\begin{keyformula}[Translation removes the constant twist; Thm.~1.2, eq.~(1.17)]\label{KF:untwist}
Write $x=e^{2\pi i\alpha}$ with $0\le\alpha<1$, put $\lambda_x:=\alpha/m$ and
$\widetilde{\mathcal H}_{d,x}(\tau,w):=\mathcal H_{d,x}(\tau,w+\lambda_x)$. Then the scalar twist
is trivialized,
\[
x^2\,e^{-4\pi i m\lambda_x}=1
\qquad(\text{since } m\lambda_x=\alpha \text{ and } x^2=e^{4\pi i\alpha}),
\]
so that
\[
\widetilde{\mathcal H}_{d,x}(\tau,w+1)=\widetilde{\mathcal H}_{d,x}(\tau,w),
\qquad
\widetilde{\mathcal H}_{d,x}(\tau,w+\tau)=q^{-m}u^{-2m}\,\widetilde{\mathcal H}_{d,x}(\tau,w).
\]
\end{keyformula}

\begin{keyformula}[The Jacobi carrier: covering identity, elliptic law, weight; Prop.~3.4, eqs.~(3.8),(3.9)]\label{KF:carrier}
With $\JJ_{m,d}$ as in Definition~\ref{def:theta}, the triple-product substitution
$\vartheta_1(\tau,z)=iq^{1/8}u^{-1/2}J(u;q)$ gives the covering formula
\[
\JJ_{m,d}(\tau,z)=q^{(m+1)/12}(q;q)_\infty^{\,2-4m}\,u^{-m}\,\frac{J(u;q)^{\,4m+1}}{J(u^d;q^d)}.
\]
The theta quasi-periodicities of Definition~\ref{def:theta} yield the index-$m$ elliptic law
\[
\JJ_{m,d}(\tau,z+1)=\JJ_{m,d}(\tau,z),
\qquad
\JJ_{m,d}(\tau,z+\tau)=q^{-m}u^{-2m}\,\JJ_{m,d}(\tau,z),
\]
and the weight exponent in the modular factor $e^{\pi i((4m+1)-d)cz^2/(c\tau+d_0)}$ is
\[
(4m+1)-d=2m,
\]
i.e.\ weight $1$, index $m$.
\end{keyformula}

\begin{keyformula}[The first new case $m=2$, $d=5$: explicit defect, residues, carrier; \S5]\label{KF:level5}
For $m=2,\ d=5,\ \omega=e^{2\pi i/5}$, with $X_j:=X_j^{(2)}(1;q)$,
\[
R_{5,1}(u;q)=(q^{-1}u^2-u^3-u^2-u-1)X_0-(u^2+u)X_1-(u^3+u^2)X_2-u^3X_3,
\]
\[
S_{5,1}(u;q):=(1-u)R_{5,1}(u;q)-5u^4\,\frac{(q^5;q^5)_\infty}{(q;q)_\infty},
\]
and Key Formula~\ref{KF:defect} specializes ($x=1$, so the $1-xu$ factor cancels) to
\[
K_{5,1}(qu;q)-K_{5,1}(u;q)=\frac{M_5(u;q)}{u(1-u^5)}\,S_{5,1}(u;q),
\qquad
D_{r,1}(q)=\frac{\omega^r}{5}\,S_{5,1}(\omega^{-r};q).
\]
The level-$5$ carrier is
\[
\JJ_{2,5}(\tau,w)=\eta(\tau)^{-6}\frac{\vartheta_1(\tau,w)^9}{\vartheta_1(5\tau,5w)}
=q^{1/4}(q;q)_\infty^{-6}u^{-2}\frac{J(u;q)^9}{J(u^5;q^5)},
\]
of weight $1$, index $2$, on $\Gamma(600)$ (here $2-4m=-6$, $4m+1=9$, $24d^2=600$).
\end{keyformula}

\begin{keyformula}[Combinatorial $2\times2$ inversion for congruence multiplicities; Cor.~A.6]\label{KF:combi}
Let $d=2m+1$. Given $U_m(N),D_m(N)\in\C$ satisfying the linear system
\[
U_m(N)=A_m(N)+(d-1)B_m(N),
\qquad
D_m(N)=A_m(N)-B_m(N),
\]
the unique solution is
\[
A_m(N)=\frac{U_m(N)+(d-1)D_m(N)}{d},
\qquad
B_m(N)=\frac{U_m(N)-D_m(N)}{d}.
\]
\end{keyformula}

\end{document}
